\textbf{Scope and chronology.} This chapter audits all 23 requested catalogue
records and all 55 supplied scan pages. The mathematical core belongs to
1976--77. One requested record does not: P000753 dates record~437 to
13 February 1978, not the catalogue's 13 February 1977. It is retained as
a chronological appendix, with corrected date and \texttt{1977-437}
occurrence identifiers. The five adjacent pages P000705--P000706,
P000752 and P000757--P000758 are accounted for without assigning their
new content to 1976--77. Plans and year headings remain visible records.
The archive list is a finding aid, not the authority for a conflicting
manuscript date \citep{mccme2026seminararchive}.

\textbf{The year's lines.} The opening talk connects real integral geometry
on \(SL(2,\mathbb R)\) with holomorphic domains and separation of discrete
series. Its strongest exact match is the 1977 Gelfand--Gindikin paper;
the notebook's projector kernel omits a nonconstant trace factor
\citep[Eq.~22 and Proposition~9]{gelfandgindikin1977complex}. Fuchs and
Kirillov revisit the Weil algebra and transgression. Arnold's discussion
of second-order equations is principally classical projective geometry:
point linearization requires the full Lie--Tresse conditions, not a
vanishing coefficient along one solution
\citep[\S\S VII--VIII]{cartan1924projective}.

Superalgebras, the Wess--Zumino multiplet, Berezin integration and Stokes
boundaries continue the preceding year's supergeometry. The distinction
between an ordinary differential form and an integrable superdensity is
essential. The Bernstein--Leites boundary theorem is a precise later
publication, while the Popov--Daikhin Einstein--Yang--Mills construction
had already appeared in 1975
\citep{bernsteinleites1977integral,popovdaikhin1975einstein}.
The latter match exposes a speaker-initial conflict that the notes
alone cannot resolve. Migdal's confinement discussion remains a physical
programme; its class E does not assert a theorem of continuum confinement.

The winter material links Euler products and affine Lie algebras with
integrable equations. Manin's filtered bimodules, spectral curves and
nodal gluing have exact counterparts in his 1978 survey, translated in
1979. Comparing formulas, rather than titles, also identifies an erroneous
flatness sign shared by the notebook and the printed survey
\citep[Ch.~III, Theorem~2.6]{manin1978algebraic}.
Duistermaat's variational index formula had circulated in May 1974;
his spectral discussion leads to the later joint paper with Kolk and
Varadarajan. Gabrielov's hypersimplicial chains and Graev's admissible
line complexes likewise belong to established programmes with earlier
publications, not newly discovered orphan problems.

The spring notes connect Poncelet closure, almost Euclidean sections,
Young diagrams and representation packets. Vershik's limit-shape
announcement is distinguished from the later Gaussian fluctuation theorem:
the latter is a generalized random field with a specified covariance,
not the notebook's unqualified white noise
\citep[\S7]{ivanovolshanski2002kerov}. Tyurin's universal assertion that
varying the complex structure can merge any two irreducible representations
fails already for distinct unitary characters; his own later survey records
the relevant uniqueness principle
\citep[Ch.~II, \S2, Exercise~10]{tyurin1978periods}. The 1978 appendix
records the BGG correspondence and distinguishes it from Beilinson's
related construction \citep{bernsteingelfandgelfand1978bundles,beilinson1978coherent}.

\begin{center}
\small
\begin{tabular}{lrr}
\hline
Historical class & Entire supplied package & 1976--77 core\\
\hline
A: published soon after & 55 & 45\\
B: later development by others & 2 & 2\\
C: partly developed & 0 & 0\\
D: not developed & 0 & 0\\
E: already known & 176 & 163\\
F: not identifiable precisely & 17 & 17\\
\hline
Total & 250 & 227\\
\hline
\end{tabular}
\end{center}

\textbf{Questions and limits.} The problem index contains 15
definite questions, including resolved questions and explicitly linked
earlier occurrences. No item meets the C/D development gate. Unspecified
range spaces, continuous-order matrices, deleted global conditions and
unqualified behaviour questions remain F; they supply no evidence of an
undeveloped theorem. The development section contains complete critical
checks and counterexamples, marked as our verifications, with no claim
of research novelty. Historical fate, literal correctness and depth of
verification are separate assessments throughout.

\textbf{Continuities.} The supplied earlier index provides particularly
useful routes through 1975 records~366 and~370 (Euler products), 374--385
and~391 (spectral curves and inverse constructions), 377 and~386 (gauge
theory), and~389 (supergeometry). Older Morse/Maslov and integral-geometric
threads provide context. The BGG question over projective space is not
merged with the earlier affine Serre--Quillen--Suslin problem. Related
vocabulary is never enough to establish identity of questions.
